% Options for packages loaded elsewhere
\PassOptionsToPackage{unicode}{hyperref}
\PassOptionsToPackage{hyphens}{url}
\PassOptionsToPackage{dvipsnames,svgnames*,x11names*}{xcolor}
%
\documentclass[
  11pt,
]{article}
\usepackage{amsmath,amssymb}
\usepackage{lmodern}
\usepackage{iftex}
\ifPDFTeX
  \usepackage[T1]{fontenc}
  \usepackage[utf8]{inputenc}
  \usepackage{textcomp} % provide euro and other symbols
\else % if luatex or xetex
  \usepackage{unicode-math}
  \defaultfontfeatures{Scale=MatchLowercase}
  \defaultfontfeatures[\rmfamily]{Ligatures=TeX,Scale=1}
  \setmainfont[]{Helvetica Neue}
  \setsansfont[]{Helvetica Neue}
  \setmathfont[]{STIX Two Math}
\fi
% Use upquote if available, for straight quotes in verbatim environments
\IfFileExists{upquote.sty}{\usepackage{upquote}}{}
\IfFileExists{microtype.sty}{% use microtype if available
  \usepackage[]{microtype}
  \UseMicrotypeSet[protrusion]{basicmath} % disable protrusion for tt fonts
}{}
\makeatletter
\@ifundefined{KOMAClassName}{% if non-KOMA class
  \IfFileExists{parskip.sty}{%
    \usepackage{parskip}
  }{% else
    \setlength{\parindent}{0pt}
    \setlength{\parskip}{6pt plus 2pt minus 1pt}}
}{% if KOMA class
  \KOMAoptions{parskip=half}}
\makeatother
\usepackage{xcolor}
\IfFileExists{xurl.sty}{\usepackage{xurl}}{} % add URL line breaks if available
\usepackage{hyperref}
\IfFileExists{bookmark.sty}{\usepackage{bookmark}}{}
\hypersetup{
  colorlinks=true,
  linkcolor={accent},
  filecolor={Maroon},
  citecolor={Blue},
  urlcolor={Blue},
  pdfcreator={LaTeX via pandoc}}
\urlstyle{same} % disable monospaced font for URLs
\usepackage[margin=18mm]{geometry}
\setlength{\emergencystretch}{3em} % prevent overfull lines
\providecommand{\tightlist}{%
  \setlength{\itemsep}{0pt}\setlength{\parskip}{0pt}}
\setcounter{secnumdepth}{3} % keep section anchors for clickable PDF TOC
\usepackage{xcolor}
\usepackage{titlesec}
\definecolor{textmain}{HTML}{1F2328}
\definecolor{muted}{HTML}{6B7280}
\definecolor{accent}{HTML}{2563EB}
\definecolor{linegray}{HTML}{D9DEE7}
\color{textmain}
\renewcommand{\familydefault}{\sfdefault}
\setlength{\parindent}{0pt}
\setlength{\parskip}{5pt}
\setlength{\emergencystretch}{3em}
\titleformat{\section}{\fontsize{16}{20}\selectfont\bfseries\color{textmain}}{}{0pt}{}
\titleformat{\subsection}{\fontsize{13}{16}\selectfont\bfseries\color{accent}}{}{0pt}{}
\titlespacing*{\section}{0pt}{18pt}{6pt}
\titlespacing*{\subsection}{0pt}{12pt}{4pt}
\ifLuaTeX
  \usepackage{selnolig}  % disable illegal ligatures
\fi

\title{Efficient Partition Trees — Calculation Notes}
\usepackage{etoolbox}
\makeatletter
\providecommand{\subtitle}[1]{% add subtitle to \maketitle
  \apptocmd{\@title}{\par {\large #1 \par}}{}{}
}
\makeatother
\subtitle{matousek-partition-tree / docs}
\author{}
\date{}

\begin{document}
\maketitle

{
\hypersetup{linkcolor=}
\setcounter{tocdepth}{3}
\tableofcontents
}
\hypertarget{efficient-partition-trees-calculation-notes}{%
\section{Efficient Partition Trees: Calculation
Notes}\label{efficient-partition-trees-calculation-notes}}

Source: Jiří Matoušek, ``Efficient Partition Trees'', Discrete \&
Computational Geometry 8:315-334, 1992.

Coverage: calculation-focused extraction. This records the quantitative
dependencies, recurrence bounds, parameter choices, and proof algebra.
It does not reproduce historical prose or bibliography details except
where they affect the bounds.

\hypertarget{global-conventions}{%
\subsection{Global Conventions}\label{global-conventions}}

\begin{itemize}
\tightlist
\item
  Dimension is fixed: d is treated as a constant.
\item
  Hidden constants may depend on d and on fixed epsilon-style constants.
\item
  P is an n-point set in Euclidean d-space.
\item
  H is usually a set of n hyperplanes.
\item
  s is the desired class size in a simplicial partition.
\item
  r is usually:
\end{itemize}

\[
r=\frac{n}{s}
\]

\begin{itemize}
\tightlist
\item
  A simplicial partition is a collection:
\end{itemize}

\[
\Pi=\lbrace (P_1,\Delta_1),\ldots,(P_m,\Delta_m) \rbrace
\]

where the sets \(P_i\) partition P, and each relatively open simplex
\(\Delta_i\) contains \(P_i\).

\begin{itemize}
\tightlist
\item
  The class-size condition is:
\end{itemize}

\[
\max_i |P_i| < 2\min_i |P_i|
\]

In the main theorem this becomes:

\[
s\le |P_i|<2s
\]

\begin{itemize}
\tightlist
\item
  A hyperplane h crosses a simplex \(\Delta\) if:
\end{itemize}

\[
h\cap \Delta\ne\emptyset
\]

and:

\[
\Delta\not\subset h
\]

\begin{itemize}
\tightlist
\item
  The crossing number of h relative to \(\Pi\) is the number of
  simplices \(\Delta_i\) crossed by h.
\item
  The crossing number of \(\Pi\) is the maximum over all hyperplanes h.
\end{itemize}

\hypertarget{section-2-cuttings}{%
\subsection{Section 2: Cuttings}\label{section-2-cuttings}}

For a set H of n hyperplanes and a cutting \(\Xi\), define:

\[
H_\Delta=\lbrace h\in H:h\text{ intersects the interior of }\Delta \rbrace
\]

A cutting \(\Xi\) is an epsilon-cutting if:

\[
|H_\Delta|\le \varepsilon n
\]

for every simplex \(\Delta\in\Xi\).

Weighted version:

\[
w(X)=\sum_{h\in X}w(h)
\]

A weighted epsilon-cutting satisfies:

\[
w(H_\Delta)\le \varepsilon w(H)
\]

Theorem 2.1 gives:

\[
\text{there exists a }(1/r)\text{-cutting of size }O(r^d)
\]

and it can be computed in:

\[
O(nr^{d-1})
\]

Weighted cuttings have the same asymptotic size and running-time bounds.

\hypertarget{randomized-cutting-construction-in-section-2}{%
\subsection{Randomized Cutting Construction in Section
2}\label{randomized-cutting-construction-in-section-2}}

The paper sketches the standard two-level cutting algorithm:

\begin{enumerate}
\def\labelenumi{\arabic{enumi}.}
\tightlist
\item
  Sample r hyperplanes from H.
\item
  Build and triangulate the arrangement of the sample.
\item
  For each sample simplex \(\Delta\), compute:
\end{enumerate}

\[
H_\Delta=\lbrace h\in H:h\text{ intersects }\Delta \rbrace
\]

\begin{enumerate}
\def\labelenumi{\arabic{enumi}.}
\setcounter{enumi}{3}
\tightlist
\item
  Define:
\end{enumerate}

\[
t_\Delta=\frac{|H_\Delta|r}{n}
\]

\begin{enumerate}
\def\labelenumi{\arabic{enumi}.}
\setcounter{enumi}{4}
\tightlist
\item
  Sample:
\end{enumerate}

\[
C t_\Delta \log t_\Delta
\]

hyperplanes from \(H_\Delta\), add the boundary hyperplanes of
\(\Delta\), triangulate inside \(\Delta\), and combine all resulting
simplices.

The resulting cutting has expected optimal size:

\[
O(r^d)
\]

and expected construction time:

\[
O(nr^{d-1})
\]

\hypertarget{section-3-partition-theorem}{%
\subsection{Section 3: Partition
Theorem}\label{section-3-partition-theorem}}

Main theorem:

Given P, n, s, and:

\[
r=\frac{n}{s}
\]

there exists a simplicial partition \(\Pi\) with:

\[
s\le |P_i|<2s
\]

and crossing number:

\[
O(r^{1-1/d})
\]

This is tight in the worst case.

\hypertarget{lemma-3.2-fixed-test-set-q}{%
\subsection{Lemma 3.2: Fixed Test Set
Q}\label{lemma-3.2-fixed-test-set-q}}

Input:

\begin{itemize}
\tightlist
\item
  P, n, s.
\item
  Q, a finite set of hyperplanes.
\end{itemize}

Output:

A simplicial partition \(\Pi\) with:

\[
s\le |P_i|<2s
\]

such that every q in Q has crossing number:

\[
O(r^{1-1/d}+\log |Q|)
\]

\hypertarget{iterative-construction}{%
\subsubsection{Iterative Construction}\label{iterative-construction}}

After i rounds:

\begin{itemize}
\tightlist
\item
  Already selected point classes:
\end{itemize}

\[
P_1,\ldots,P_i
\]

\begin{itemize}
\tightlist
\item
  Remaining points:
\end{itemize}

\[
P'_i=P\setminus(P_1\cup\cdots\cup P_i)
\]

\begin{itemize}
\tightlist
\item
  Remaining count:
\end{itemize}

\[
n_i=|P'_i|
\]

If:

\[
n_i<2s
\]

then set the last class to all remaining points and terminate.

For q in Q, let:

\[
x_i(q)=\#\lbrace \Delta_1,\ldots,\Delta_i\text{ crossed by }q \rbrace
\]

Define weights:

\[
w_i(q)=2^{x_i(q)}
\]

Thus each time q crosses a newly selected simplex, its weight doubles.

Let total weight be:

\[
w_i(Q)=\sum_{q\in Q}w_i(q)
\]

Choose \(t_i\) as large as possible such that there exists a weighted
\((1/t_i)\)-cutting for \((Q,w_i)\) whose simplices have at most:

\[
\frac{n_i}{s}
\]

faces of all dimensions in total.

By Theorem 2.1, since a cutting with parameter \(t_i\) has \(O(t_i^d)\)
faces, \(t_i\) can satisfy:

\[
t_i>c\left(\frac{n_i}{s}\right)^{1/d}
\]

for some positive constant c.

Since the cutting has at most:

\[
\frac{n_i}{s}
\]

faces and there are \(n_i\) remaining points, the pigeonhole principle
gives a relatively open face containing at least s remaining points.

Choose such a face as:

\[
\Delta_{i+1}
\]

Pick exactly s points inside it:

\[
P_{i+1}\subset P'_i\cap\Delta_{i+1}
\]

Then continue.

\hypertarget{weight-growth}{%
\subsubsection{Weight Growth}\label{weight-growth}}

Let:

\[
Q_{i+1}=\lbrace q\in Q:q\text{ crosses }\Delta_{i+1} \rbrace
\]

For q in \(Q_{i+1}\), the weight doubles. For all other q, the weight
stays the same.

Therefore:

\[
w_{i+1}(Q) =w_i(Q)-w_i(Q_{i+1})+2w_i(Q_{i+1})
\]

so:

\[
w_{i+1}(Q) =w_i(Q)+w_i(Q_{i+1})
\]

and:

\[
w_{i+1}(Q) =w_i(Q)\left(1+\frac{w_i(Q_{i+1})}{w_i(Q)}\right)
\]

Because \(\Delta_{i+1}\) is a face of a weighted \((1/t_i)\)-cutting:

\[
w_i(Q_{i+1})\le \frac{w_i(Q)}{t_i}
\]

For full-dimensional faces this is exactly the cutting guarantee. For
lower-dimensional faces, a hyperplane crossing the face also intersects
the interior of an adjacent full-dimensional simplex, so the same bound
holds up to a constant absorbed in c.

Thus:

\[
w_{i+1}(Q) \le w_i(Q)\left(1+\frac{1}{t_i}\right)
\]

Using:

\[
t_i>c\left(\frac{n_i}{s}\right)^{1/d}
\]

we get:

\[
\frac{1}{t_i} < \frac{1}{c}\left(\frac{s}{n_i}\right)^{1/d}
\]

In nonterminal rounds exactly s points are removed, so:

\[
n_i=n-is
\]

and:

\[
r=\frac{n}{s}
\]

The number of nonterminal rounds is:

\[
m=\lfloor r\rfloor
\]

Hence:

\[
n_i=s(r-i)
\]

so:

\[
\left(\frac{s}{n_i}\right)^{1/d} = \frac{1}{(r-i)^{1/d}}
\]

Iterating the weight inequality:

\[
w_m(Q) \le |Q|\prod_{i=0}^{m-1}\left(1+\frac{1}{c(r-i)^{1/d}}\right)
\]

Taking logarithms and using:

\[
\ln(1+x)\le x
\]

gives:

\[
\log w_m(Q) \le \log |Q| + \frac{1}{c}\sum_{i=0}^{m-1}\frac{1}{(r-i)^{1/d}}
\]

Rewrite the sum by \(j=r-i\):

\[
\sum_{i=0}^{m-1}\frac{1}{(r-i)^{1/d}} \le \sum_{j=1}^{r}\frac{1}{j^{1/d}}
\]

Bound by integral:

\[
\sum_{j=1}^{r}j^{-1/d} = O(r^{1-1/d})
\]

Therefore:

\[
\log w_m(Q) = O(\log |Q|+r^{1-1/d})
\]

If a hyperplane q has final crossing number x, then:

\[
w_m(q)=2^x
\]

Since:

\[
w_m(q)\le w_m(Q)
\]

we get:

\[
x\le \log_2 w_m(Q) = O(\log |Q|+r^{1-1/d})
\]

This proves Lemma 3.2.

\hypertarget{lemma-3.3-test-set-lemma}{%
\subsection{Lemma 3.3: Test Set Lemma}\label{lemma-3.3-test-set-lemma}}

Purpose: reduce all possible hyperplanes to a finite test set Q.

Input:

\begin{itemize}
\tightlist
\item
  P, n.
\item
  Parameter r.
\end{itemize}

Output:

A set Q of at most r hyperplanes such that if a simplicial partition has
all classes of size at least s, and if:

\[
x_0=\max_{q\in Q}\mathrm{cr}_\Pi(q)
\]

then every hyperplane h has crossing number:

\[
\mathrm{cr}_\Pi(h) \le (d+1)x_0 + O\left(\frac{n}{s r^{1/d}}\right)
\]

\hypertarget{construction-of-q}{%
\subsubsection{Construction of Q}\label{construction-of-q}}

Dualize P:

\[
H=\mathcal{D}(P)
\]

This is a set of n hyperplanes.

Choose a cutting for H with cutting parameter:

\[
t=\Omega(r^{1/d})
\]

The cutting has \(O(t^d)\) vertices. Choose the constant so that the
total number of vertices is at most r.

Let V be the set of vertices of the cutting. Define:

\[
Q=\mathcal{D}(V)
\]

Then:

\[
|Q|\le r
\]

\hypertarget{why-q-controls-every-h}{%
\subsubsection{Why Q Controls Every h}\label{why-q-controls-every-h}}

Take an arbitrary hyperplane h.

Its dual point:

\[
\mathcal{D}(h)
\]

lies in some cutting simplex A.

Let G be the set of vertices of A. A d-dimensional simplex has \(d+1\)
vertices.

The dual hyperplanes:

\[
\mathcal{D}(G)
\]

belong to Q.

Each of these \(d+1\) test hyperplanes crosses at most \(x_0\) simplices
of \(\Pi\), so simplices crossed by at least one of them are at most:

\[
(d+1)x_0
\]

Now consider simplices crossed by h but by no hyperplane in
\(\mathcal{D}(G)\).

Such simplices must be completely contained in the zone of h in the
arrangement of \(\mathcal{D}(G)\).

By duality, any point p of P lying in the interior of that zone has dual
hyperplane \(\mathcal{D}(p)\) intersecting the interior of A.

Because A is a \((1/t)\)-cutting cell and:

\[
t=\Omega(r^{1/d})
\]

the number of dual hyperplanes crossing A is:

\[
O\left(\frac{n}{r^{1/d}}\right)
\]

Thus the number of original points in those exceptional simplices is:

\[
O\left(\frac{n}{r^{1/d}}\right)
\]

Each partition class has at least s points, so the number of exceptional
simplices is:

\[
O\left(\frac{n}{s r^{1/d}}\right)
\]

Therefore:

\[
\mathrm{cr}_\Pi(h) \le (d+1)x_0 + O\left(\frac{n}{s r^{1/d}}\right)
\]

This is the Test Set Lemma.

\hypertarget{proof-of-theorem-3.1}{%
\subsection{Proof of Theorem 3.1}\label{proof-of-theorem-3.1}}

Use Lemma 3.3 to get:

\[
|Q|\le r
\]

Then use Lemma 3.2 on this Q.

For every q in Q:

\[
\mathrm{cr}_\Pi(q) = O(r^{1-1/d}+\log |Q|)
\]

Since:

\[
|Q|\le r
\]

we have:

\[
\log |Q|\le \log r
\]

and for fixed d at the relevant scale:

\[
\log r=O(r^{1-1/d})
\]

so:

\[
x_0=O(r^{1-1/d})
\]

Now apply Lemma 3.3:

\[
\mathrm{cr}_\Pi(h) \le (d+1)O(r^{1-1/d}) + O\left(\frac{n}{s r^{1/d}}\right)
\]

Since:

\[
r=\frac{n}{s}
\]

we have:

\[
\frac{n}{s r^{1/d}} = \frac{r}{r^{1/d}} = r^{1-1/d}
\]

Therefore:

\[
\mathrm{cr}_\Pi(h) = O(r^{1-1/d})
\]

for every hyperplane h.

\hypertarget{tightness-argument-in-theorem-3.1}{%
\subsection{Tightness Argument in Theorem
3.1}\label{tightness-argument-in-theorem-3.1}}

Choose a set Q of hyperplanes in general position whose arrangement has
at least:

\[
\left\lceil\frac{n}{s-1}\right\rceil
\]

distinct cells.

It suffices to take:

\[
O(r^{1/d})
\]

hyperplanes, because an arrangement of M hyperplanes in fixed dimension
d has:

\[
\Theta(M^d)
\]

cells.

Place n points into clusters of at most \(s-1\) points each, with
clusters in distinct cells of the arrangement of Q.

In any simplicial partition with classes of size at least s, no class
can be contained in one such cluster. Therefore every simplex must cross
at least one hyperplane of Q.

There are about:

\[
\frac{n}{s}=r
\]

simplices/classes.

The total number of crossings over all hyperplanes in Q is at least:

\[
\Omega(r)
\]

Since:

\[
|Q|=O(r^{1/d})
\]

the average crossing number of a hyperplane in Q is:

\[
\Omega\left(\frac{r}{r^{1/d}}\right) = \Omega(r^{1-1/d})
\]

Thus the upper bound is asymptotically tight.

\hypertarget{lemma-3.4-constant-r-construction-time}{%
\subsection{Lemma 3.4: Constant r Construction
Time}\label{lemma-3.4-constant-r-construction-time}}

If r is bounded by a constant:

\begin{itemize}
\tightlist
\item
  Test Set Lemma needs a cutting with constant parameter, computable in:
\end{itemize}

\[
O(n)
\]

\begin{itemize}
\tightlist
\item
  Q has constant size.
\item
  Lemma 3.2 makes only constant many steps and handles constant many
  simplices.
\end{itemize}

The only issue is large weights:

\[
w_i(q)=2^{x_i(q)}
\]

For weighted cutting computation, it is enough to know relative weights
within a constant factor, because the weighted algorithm replaces h by
approximately:

\[
m(h)=\left\lfloor\frac{|H|w(h)}{w(H)}\right\rfloor+1
\]

copies.

Thus ordinary integer arithmetic with constant relative accuracy is
enough.

Conclusion:

\[
O(n)
\]

time for constant r.

\hypertarget{corollary-3.5-simple-algorithmic-partition}{%
\subsection{Corollary 3.5: Simple Algorithmic
Partition}\label{corollary-3.5-simple-algorithmic-partition}}

Goal:

Construct a partition with class size s and crossing number:

\[
O(r^{1-1/d+\delta})
\]

in:

\[
O(n\log r)
\]

for fixed positive \(\delta\).

\hypertarget{composition-calculation}{%
\subsubsection{Composition Calculation}\label{composition-calculation}}

First build a partition with:

\[
r_1=\frac{n}{s_1}
\]

and crossing number:

\[
C r_1^{1-1/d}
\]

Then refine each class with:

\[
r_2=\frac{2s_1}{s_2}
\]

and crossing number:

\[
C r_2^{1-1/d}
\]

The combined crossing number is bounded by a constant times:

\[
C^2 r_1^{1-1/d}r_2^{1-1/d} = O((r_1r_2)^{1-1/d})
\]

The paper writes the resulting multiplicative loss as:

\[
2C^2 r^{1-1/d}
\]

where:

\[
r=\frac{n}{s_2}
\]

Repeat with a fixed constant reduction factor \(r_0\).

Number of iterations:

\[
D=\frac{\log r}{\log r_0}
\]

Each iteration loses a constant factor, so the total loss is:

\[
(2C)^D = r^{\log(2C)/\log r_0}
\]

Choose \(r_0\) large enough so:

\[
\frac{\log(2C)}{\log r_0}< \delta
\]

Then final crossing number:

\[
O(r^{1-1/d+\delta})
\]

Total time over levels:

\[
O(n\log r)
\]

\hypertarget{section-4-epsilon-approximations}{%
\subsection{Section 4:
Epsilon-Approximations}\label{section-4-epsilon-approximations}}

An epsilon-approximation A for P with respect to simplices satisfies:

\[
\left| \frac{|A\cap\sigma|}{|A|} - \frac{|P\cap\sigma|}{|P|} \right| < \varepsilon
\]

Weighted form:

\[
\left| \frac{w(A\cap\sigma)}{w(A)} - \frac{|P\cap\sigma|}{|P|} \right| < \varepsilon
\]

\hypertarget{observation-4.1}{%
\subsection{Observation 4.1}\label{observation-4.1}}

Given a simplicial partition:

\[
\Pi=\lbrace (P_1,\Delta_1),\ldots,(P_m,\Delta_m) \rbrace
\]

with:

\[
|P_i|\le s
\]

and crossing number x.

Pick one representative point:

\[
a_i\in P_i
\]

with weight:

\[
w(a_i)=|P_i|
\]

Let:

\[
A=\lbrace a_1,\ldots,a_m \rbrace
\]

For a query simplex \(\sigma\), the representative count differs from
the true count only for partition simplices crossed by the boundary of
\(\sigma\).

A d-dimensional simplex has \(d+1\) bounding hyperplanes. Each boundary
hyperplane crosses at most x partition simplices.

So at most:

\[
(d+1)x
\]

classes can contribute error.

Each class has at most s points, so absolute error is at most:

\[
(d+1)xs
\]

Relative error is:

\[
\varepsilon=\frac{(d+1)xs}{n}
\]

Thus:

\[
(A,w)
\]

is an epsilon-approximation with:

\[
\varepsilon=\frac{(d+1)xs}{n}
\]

\hypertarget{theorem-4.2}{%
\subsection{Theorem 4.2}\label{theorem-4.2}}

Using Corollary 3.5 and Observation 4.1, compute a
\((1/t)\)-approximation of size:

\[
O(t^{d+\delta})
\]

in:

\[
O(n\log t)
\]

The parameter is chosen so that the partition crossing error satisfies:

\[
\frac{(d+1)xs}{n}\le \frac{1}{t}
\]

with:

\[
x=O(r^{1-1/d+\delta})
\]

and:

\[
r=\frac{n}{s}
\]

This yields representative set size:

\[
O(r)=O(t^{d+\delta})
\]

after absorbing constant-dimension exponent slack.

\hypertarget{lemma-4.3-approximation-implies-cutting}{%
\subsection{Lemma 4.3: Approximation Implies
Cutting}\label{lemma-4.3-approximation-implies-cutting}}

For hyperplanes H, let A be an epsilon-approximation for H with respect
to segments.

If \(\Xi\) is an epsilon-prime cutting for weighted A, then \(\Xi\) is
a:

\[
d(\varepsilon+\varepsilon')
\]

cutting for H.

The factor d comes from reducing simplex-cell crossing control to
segment-edge crossing control in fixed dimension.

\hypertarget{proposition-4.4}{%
\subsection{Proposition 4.4}\label{proposition-4.4}}

Goal: compute a \((1/r)\)-cutting of size:

\[
O(r^d)
\]

for n hyperplanes.

Step 1: compute a \((1/(2dr))\)-approximation of H of size:

\[
O(r^{d+\delta})
\]

in time:

\[
O(n\log r)
\]

Step 2: compute a \((1/(2dr))\)-cutting for the approximation.

The approximation size is:

\[
N_A=O(r^{d+\delta})
\]

Using Theorem 2.1:

\[
O(N_A r^{d-1}) = O(r^{d+\delta}r^{d-1}) = O(r^{2d+\delta-1})
\]

Total:

\[
O(n\log r+r^{2d+\delta-1})
\]

If:

\[
r< n^{1/(2d+\delta-1)}
\]

then:

\[
r^{2d+\delta-1}< n
\]

so total time is:

\[
O(n\log r)
\]

\hypertarget{lemma-4.5-auxiliary-range-counting}{%
\subsection{Lemma 4.5: Auxiliary Range
Counting}\label{lemma-4.5-auxiliary-range-counting}}

Given:

\begin{itemize}
\tightlist
\item
  fixed constant C,
\item
  point set P of size n,
\item
  parameter r satisfying:
\end{itemize}

\[
r< n^\alpha
\]

for sufficiently small \(\alpha=\alpha(C)\).

Build in:

\[
O(n\log r)
\]

a structure answering simplex counting queries in:

\[
O\left(\frac{n}{r^C}\right)
\]

and supporting deletions in:

\[
O(1)
\]

per deletion.

Construction:

Choose t as a sufficiently large power of r.

Build a simplicial partition with at most t classes and crossing number:

\[
O(t^{1-1/d+\delta})
\]

Each class has size about:

\[
\frac{n}{t}
\]

For a query simplex, directly inspect classes crossed by the boundary,
and add stored counts for fully contained classes.

Query time:

\[
O\left(t+\frac{n}{t}t^{1-1/d+\delta}\right)
\]

Simplify:

\[
O\left(t+n t^{-1/d+\delta}\right)
\]

Choose t as a sufficiently large power of r and choose \(\alpha\) small
enough so:

\[
O\left(t+n t^{-1/d+\delta}\right) \le O\left(\frac{n}{r^C}\right)
\]

\hypertarget{lemma-4.6-fast-partition-for-small-r}{%
\subsection{Lemma 4.6: Fast Partition for Small
r}\label{lemma-4.6-fast-partition-for-small-r}}

If:

\[
r< n^\alpha
\]

for sufficiently small fixed \(\alpha>0\), then the exact Partition
Theorem partition can be constructed in:

\[
O(n\log r)
\]

Key costs:

\begin{enumerate}
\def\labelenumi{\arabic{enumi}.}
\tightlist
\item
  Test-set cutting with:
\end{enumerate}

\[
t=\Omega(r^{1/d})
\]

computed using Proposition 4.4 in:

\[
O(n\log r)
\]

\begin{enumerate}
\def\labelenumi{\arabic{enumi}.}
\setcounter{enumi}{1}
\tightlist
\item
  In Lemma 3.2, selecting a face containing enough remaining points
  requires counting points in cutting faces across rounds.
\end{enumerate}

There are:

\[
O(r^2)
\]

range counting queries total in the argument as stated in the paper.

Using Lemma 4.5, all dynamic range-counting and reporting operations are
covered within:

\[
O(n\log r)
\]

for sufficiently small \(\alpha\).

\hypertarget{theorem-4.7}{%
\subsection{Theorem 4.7}\label{theorem-4.7}}

Three algorithmic partition bounds.

\hypertarget{part-1}{%
\subsubsection{Part 1}\label{part-1}}

If:

\[
s\ge n^\delta
\]

then:

\[
r=\frac{n}{s}\le n^{1-\delta}
\]

The construction can be iterated a constant number of times using Lemma
4.6.

Result:

\begin{itemize}
\tightlist
\item
  class sizes:
\end{itemize}

\[
s\le |P_i|<2s
\]

\begin{itemize}
\tightlist
\item
  crossing number:
\end{itemize}

\[
O(r^{1-1/d})
\]

\begin{itemize}
\tightlist
\item
  time:
\end{itemize}

\[
O(n\log r)
\]

\hypertarget{part-2}{%
\subsubsection{Part 2}\label{part-2}}

For arbitrary s, iterate the partition construction at most:

\[
O(\log\log n)
\]

times.

Each iteration loses a constant factor in crossing number.

Total multiplicative loss:

\[
2^{O(\log\log n)} = (\log n)^{O(1)}
\]

Result:

\[
O(r^{1-1/d}(\log r)^{O(1)})
\]

crossing number in:

\[
O(n\log r)
\]

time.

\hypertarget{part-3}{%
\subsubsection{Part 3}\label{part-3}}

For arbitrary s, exact optimal crossing:

\[
O(r^{1-1/d})
\]

can be constructed in:

\[
O(n^{1+\delta})
\]

Idea:

\begin{itemize}
\tightlist
\item
  Iterate fast construction until residual class sizes are small:
\end{itemize}

\[
O(n^{\delta'})
\]

\begin{itemize}
\tightlist
\item
  On small classes, use the polynomial-time exact construction from
  Theorem 3.1.
\item
  Choose \(\delta'\) small enough so total time is:
\end{itemize}

\[
O(n^{1+\delta})
\]

\hypertarget{theorem-4.8-improved-cutting-computation}{%
\subsection{Theorem 4.8: Improved Cutting
Computation}\label{theorem-4.8-improved-cutting-computation}}

For:

\[
r\le n^{1-\delta}
\]

a \((1/r)\)-cutting of size:

\[
O(r^d)
\]

can be computed in:

\[
O(n\log r+\sqrt{n}r^{d-1/2})
\]

Calculation:

Choose intermediate parameter:

\[
r_1= \left\lfloor \frac{r^{1+1/(2(d-1))}}{n^{1/(2(d-1))}} \right\rfloor
\]

Compute a \((1/r_1)\)-cutting of H and conflict lists \(H_\Delta\).

Then for each cell \(\Delta\), compute a \((r_1/r)\)-cutting for:

\[
H_\Delta\cup B_\Delta
\]

where \(B_\Delta\) are boundary hyperplanes of \(\Delta\).

Balancing the top-level and refinement costs yields:

\[
O(n\log r+\sqrt{n}r^{d-1/2})
\]

\hypertarget{theorem-5.1-basic-partition-tree}{%
\subsection{Theorem 5.1: Basic Partition
Tree}\label{theorem-5.1-basic-partition-tree}}

Build a recursive partition tree.

At a node v with m points, choose:

\[
s=\lfloor m^{1/d}\rfloor
\]

Then:

\[
r=\frac{m}{s}=m^{1-1/d}
\]

Class size after one level:

\[
m^{1/d}
\]

Thus node sizes decrease by repeated d-th roots:

\[
n,\quad n^{1/d},\quad n^{1/d^2},\quad\ldots
\]

Tree depth:

\[
O(\log\log n)
\]

At a node with m points, the crossing number of the partition is:

\[
O(r^{1-1/d}(\log r)^{O(1)})
\]

Using:

\[
r=m^{1-1/d}
\]

this is:

\[
O\left(m^{(1-1/d)^2}(\log m)^{O(1)}\right)
\]

But the recurrence in the paper is stated abstractly as:

\[
T(n)\le O(r)+O(r^{1-1/d})T(s)
\]

with:

\[
r=\frac{n}{s}
\]

and:

\[
s=n^{1/d}
\]

So:

\[
r=n^{1-1/d}
\]

and:

\[
T(n) \le O(n^{1-1/d}) + O(n^{(1-1/d)^2})T(n^{1/d})
\]

The solution is:

\[
T(n)=O(n^{1-1/d}(\log n)^{O(1)})
\]

Space:

An m-point node stores:

\[
O(r)=O(m^{1-1/d})
\]

simplices and cumulative weights.

Summing over levels gives:

\[
O(n)
\]

space.

Preprocessing:

An m-point node costs:

\[
O(m\log m)
\]

The node sizes decrease double-exponentially with depth, so total
preprocessing is:

\[
O(n\log n)
\]

\hypertarget{corollary-5.2-spacequery-tradeoff}{%
\subsection{Corollary 5.2: Space/Query
Tradeoff}\label{corollary-5.2-spacequery-tradeoff}}

Combine partition tree with a fast simplex range searching structure.

For storage parameter m satisfying:

\[
n< m< n^d
\]

preprocessing and space:

\[
O(m^{1+\delta})
\]

query time:

\[
O\left(\frac{n}{m^{1/d}}\log^{d+1}n\right)
\]

The paper writes the same tradeoff in asymptotic form with
fixed-dimension constants.

\hypertarget{section-6-reducing-query-time}{%
\subsection{Section 6: Reducing Query
Time}\label{section-6-reducing-query-time}}

Theorem 6.1 gives linear space but improved query factors, with worse
preprocessing.

General semigroup:

\[
O(n^{1-1/d}(\log\log n)^{O(1)})
\]

with:

\[
O(n^{1+\delta})
\]

preprocessing and:

\[
O(n)
\]

space.

If weights can be subtracted:

\[
O(n^{1-1/d}2^{O(\log^\ast  n)})
\]

For reporting:

\[
O(n^{1-1/d}2^{O(\log^\ast  n)}+k)
\]

\hypertarget{why-theorem-5.1-has-a-log-factor}{%
\subsection{Why Theorem 5.1 Has a Log
Factor}\label{why-theorem-5.1-has-a-log-factor}}

At each level, the partition crossing bound loses a constant factor.

With:

\[
O(\log\log n)
\]

levels, this becomes:

\[
(\log n)^{O(1)}
\]

To reduce this, use fewer levels and larger one-level partitions. But
then the query algorithm can no longer inspect all simplices in a node
directly.

Two auxiliary problems are introduced:

\begin{itemize}
\tightlist
\item
  A: report simplices crossed by the boundary of the query simplex.
\item
  B: sum weights of simplices fully contained in the query simplex.
\end{itemize}

\hypertarget{lemma-6.2-k-tuple-containment-structure}{%
\subsection{Lemma 6.2: k-Tuple Containment
Structure}\label{lemma-6.2-k-tuple-containment-structure}}

Input:

Set A of n ordered k-tuples of points, each with semigroup weight.

Query:

Sum weights of k-tuples completely contained in a query simplex.

Output:

Space and preprocessing:

\[
O(n(\log n)^{C_k})
\]

Query:

\[
O(n^{1-1/d}\exp(C'_k\sqrt{\log n}))
\]

Construction:

Induct on k.

For k equals 1, use Theorem 5.1.

For k greater than 1:

\begin{itemize}
\tightlist
\item
  Let F be the set of first components.
\item
  Build a partition tree on F.
\item
  At an m-point node choose:
\end{itemize}

\[
r=\frac{m}{s}=\exp(\sqrt{\log m})
\]

For each class \(P_i\), store a structure for the remaining
\((k-1)\)-tuples.

Query recurrence:

\[
T_k(n) \le O(r) + O(r^{1-1/d})T_k(s) + O(r)T_{k-1}\left(\frac{2n}{r}\right)
\]

with:

\[
r=\exp(\sqrt{\log n})
\]

Using induction:

\[
T_{k-1}(n) = O(n^{1-1/d}\exp(C'_{k-1}\sqrt{\log n}))
\]

The recurrence gives:

\[
T_k(n) = O(n^{1-1/d}\exp(C'_k\sqrt{\log n}))
\]

\hypertarget{lemma-6.3-reporting-segments-crossing-a-hyperplane}{%
\subsection{Lemma 6.3: Reporting Segments Crossing a
Hyperplane}\label{lemma-6.3-reporting-segments-crossing-a-hyperplane}}

Input:

n segments in d-space.

Output:

Report all segments crossing a query hyperplane.

Space and preprocessing:

\[
O(n(\log n)^{O(1)})
\]

Query:

\[
O(n^{1-1/d}(\log n)^{O(1)}+k)
\]

Construction:

\begin{itemize}
\tightlist
\item
  Build a partition tree on left endpoints.
\item
  At an m-point node choose:
\end{itemize}

\[
s=\lfloor m^{2/3}\rfloor
\]

so:

\[
r=m^{1/3}
\]

\begin{itemize}
\tightlist
\item
  For each class, store a halfspace range reporting structure on
  corresponding right endpoints.
\end{itemize}

At an m-point node, there are at most:

\[
r=m^{1/3}
\]

secondary halfspace queries, each on:

\[
O(s)=O(m^{2/3})
\]

points.

The additive output terms sum to:

\[
O(r^{1-1/d})
\]

The overhead terms sum to:

\[
O\left(m^{1/3}(m^{2/3})^{1-1/\lfloor d/2\rfloor}(\log n)^{O(1)}\right)
\]

The paper states this is bounded by:

\[
O(m^{1-1/d})
\]

for the needed fixed-dimensional comparison.

Depth:

\[
O(\log\log n)
\]

Final query:

\[
O(n^{1-1/d}(\log n)^{O(1)}+k)
\]

\hypertarget{theorem-6.1-final-calculation}{%
\subsection{Theorem 6.1 Final
Calculation}\label{theorem-6.1-final-calculation}}

Equip each node of the main partition tree with secondary structures for
Problems A and B.

At an m-point node, choose class size s so that:

\begin{itemize}
\tightlist
\item
  work in the node is:
\end{itemize}

\[
O(m^{1-1/d})
\]

\begin{itemize}
\tightlist
\item
  total secondary space/preprocessing at the node is:
\end{itemize}

\[
O\left(\frac{m}{\log m}\right)
\]

Then total space remains:

\[
O(n)
\]

and preprocessing becomes:

\[
O(n^{1+\delta})
\]

If node work is:

\[
O(m^{1-1/d})
\]

and tree depth is D, query time becomes:

\[
O(n^{1-1/d}2^{O(D)})
\]

For general semigroup weights, the auxiliary structures allow:

\[
s=2^{C\sqrt{\log m}}
\]

Passing one level down reduces:

\[
\log m
\]

to a constant multiple of:

\[
\sqrt{\log m}
\]

Thus depth:

\[
D=O(\log\log\log n)
\]

So:

\[
2^{O(D)}=(\log\log n)^{O(1)}
\]

and query time:

\[
O(n^{1-1/d}(\log\log n)^{O(1)})
\]

If weights can be subtracted, or for reporting, Problem B is faster and
we may choose:

\[
s=(\log m)^{O(1)}
\]

Then one level reduces node size from m to:

\[
(\log m)^{O(1)}
\]

Depth:

\[
D=O(\log^\ast  n)
\]

Thus query time:

\[
O(n^{1-1/d}2^{O(\log^\ast  n)})
\]

For reporting, add output size:

\[
O(n^{1-1/d}2^{O(\log^\ast  n)}+k)
\]

\hypertarget{section-7-dynamization}{%
\subsection{Section 7: Dynamization}\label{section-7-dynamization}}

Theorem 7.1:

The simplex range searching structure of Theorem 5.1 can be maintained
under insertions and deletions with:

\[
O(\log n)
\]

amortized deletion time, and:

\[
O(\log^2 n)
\]

amortized insertion time.

\hypertarget{deletions}{%
\subsubsection{Deletions}\label{deletions}}

If semigroup weights can be subtracted, update cumulative weights
directly along the search path.

Without subtraction:

\begin{itemize}
\tightlist
\item
  Store child weights in a balanced binary tree at each node.
\item
  Updating one child sum costs:
\end{itemize}

\[
O(\log r)
\]

at an m-point node.

Since node sizes shrink doubly exponentially down the partition tree,
summing over a root-to-leaf path gives:

\[
O(\log n)
\]

Assign null weights to simulate deletion.

After:

\[
\frac{n}{2}
\]

deletions, rebuild from scratch to maintain query bounds.

\hypertarget{insertions}{%
\subsubsection{Insertions}\label{insertions}}

The static structure is buildable in:

\[
O(n\log n)
\]

The range searching problem is decomposable:

\[
P=P_1\cup P_2
\]

and the answer on P combines answers on \(P_1\) and \(P_2\).

Using the Bentley/Overmars dynamization method:

\[
O(\log^2 n)
\]

amortized insertion time.

\end{document}
